\textbf{\textit{Toutes les réponses de cet exercice devront être justifiées.}}

\bigskip
Un musée propose trois formules de visites guidées pour des classes durant l’année scolaire. 

\medskip
\fbox{
\begin{minipage}{\linewidth}
Formule A : 45~\euro{} par visite de classe. 

Formule B : Abonnement annuel de 90~\euro{} par école auquel s'ajoute un montant de 25~\euro{}~50 par visite de classe de l'école. 

Formule C : Abonnement annuel d'un montant de 300~\euro{} qui permet autant de visites que le souhaite l'école.	
\end{minipage}
}

\medskip
\begin{enumerate}
	\item Une école est composée de quatre classes.
	\begin{enumerate}
		\item Si chaque classe effectue une visite, quelle formule est la plus avantageuse ? 
		\item Si chaque classe effectue deux visites, quelle formule est la plus avantageuse ?
	\end{enumerate}
	\item À partir de combien de visites de classe la formule C est-elle plus économique que la formule~B ?
	\item En vous aidant de la représentation ci-dessous, déterminer graphiquement le nombre de visites de classe à partir duquel la formule B est plus économique que la formule A.
	
	\begin{tikzpicture}
\begin{axis}[
    axis lines = left,
    xmin=0, xmax=12,
    ymin=0, ymax=320,
    xtick={0,2,...,12},
    ytick={0,50,...,300},
    grid=major,
    width=16cm,
    height=16cm
    ]

\addplot [
    domain=0:12, 
    samples=100, 
    color=blue,
    thick
    ]
    {45*x};

\addplot [
    domain=0:12, 
    samples=100, 
    color=black,
    thick
    ]
    {90+25.5*x};

\addplot [
    domain=0:12, 
    samples=100, 
    color=red,
    thick
    ]
    {300};

\end{axis}
\end{tikzpicture}

	\item L'équipe pédagogique de l'école de quatre classes décide d'organiser deux visites par classe. Elle choisit la formule la plus avantageuse et bénéficie d’une subvention de 15~\% du prix à payer de la mairie. Quel montant l’école doit-elle prévoir ?
\end{enumerate}